\section{The moderate-row seed}\label{sec:moderate}

\begin{theorem}\label{thm:moderate}
There are absolute constants $A,C_m,\eps_0>0$ such that the following holds.
Let $U\in\R^{n\times d}$ be Parseval with $2d\le n$, $d\ge A\log(2n)$, and
$|\norm{u_i}^2-a|\le\eps a$ for all $i$, where $0<\eps\le\eps_0$. Then there is
an ENP frame $\widehat U$ with $\fro{\widehat U-U}^2\le C_m\eps d$.
\end{theorem}

Throughout this section $P=UU^T$, $p=\diag P$, $D_p=\Diag(p)$, $L=L_P$, and
$\eps_0\le\frac12$, so $\frac a2\le p_i\le\frac{3a}2\le\frac34$. For a symmetric
$M\in\R^{n\times n}$, $\Lap(M)=\sum_{i<j}M_{ij}^2(e_i-e_j)(e_i-e_j)^T$ is the
Laplacian with weights $M_{ij}^2$; thus $L=\Lap(P)$,
$x^T\Lap(M)x=\frac12\fro{[\Diag x,M]}^2$, and
\begin{equation}\label{eq:lap-basic}
 \Lap(M)\preceq2\max_i\norm{M_i}^2\,I,\qquad
 \Lap(M+M')\succeq\tfrac12\Lap(M)-\Lap(M'),
\end{equation}
by $(x_i-x_j)^2\le2x_i^2+2x_j^2$ and $(v+w)^2\ge\frac12v^2-w^2$. We write
$\mathsf K=nI-\one\one^T$ for the Laplacian of the complete graph.

\paragraph{Scales.} Put $\sigma=at^2$, where $t\asymp\sqrt\eps$ is the size of
the perturbation. The perturbation creates edges of weight $\gtrsim\sigma/n$,
the resulting barrier constant is $O(1/\sigma)$, and so all diagonal errors must
be at most a small multiple of $\sigma$. The table lists the errors and how
each is controlled ($\rho=D^2t^2$ is the filter level, $D$ a large constant).
\begin{center}
\small
\begin{tabular}{@{}lll@{}}
\toprule
diagonal error & size$/\sigma$ & reason\\
\midrule
input & $\eps/t^2$ & $t^2=K_0\eps$, $K_0$ large\\
linear term $2t\mathcal AZ$ & $6D\sqrt{\log(2n)/d}$ & filter (\cref{lem:filter}), $d\ge A\log 2n$\\
quadratic fluctuation & $\eta'$ & concentration, $d\ge A\log 2n$\\
filter bias $t^2b_*$ & $0$ & cancelled by the drift (\cref{lem:mean,lem:retraction})\\
drift cross terms & $40/D+1/(2D^2)$ & $D$ large\\
base bias $t^2b_0$ & $12\eps/d$ & \cref{lem:mean}\\
retraction remainder & $C_Et$ & $t$ small\\
\bottomrule
\end{tabular}
\end{center}

\subsection{The horizontal tangent space and the filtered covariance}

Fix an isometry $V\in\R^{n\times(n-d)}$ onto $(\im U)^\perp$, and let
$\mathcal Z=\R^{(n-d)\times d}$ with the Frobenius inner product. For
$Z\in\mathcal Z$ the matrix $VZ$ is horizontal ($U^TVZ=0$), and we put
\[
 Y_Z=VZU^T+UZ^TV^T,\qquad \norm{(Y_Z)_i}^2=\norm{Z^Tv_i}^2+\norm{Zu_i}^2\le\opn Z^2,
 \qquad \fro{Y_Z}^2=2\fro Z^2 .
\]
Define
\[
 \mathcal A:\mathcal Z\to\R^n,\quad (\mathcal AZ)_i=v_i^TZu_i=(VZU^T)_{ii},\qquad
 \mathcal A^*x=N_x:=V^T\Diag(x)U .
\]
Then $\mathcal A\mathcal A^*=L$, i.e.\ $\fro{N_x}^2=x^TLx$ (\cref{lem:energy}). Let
\[
 \bar{\mathcal A}=D_p^{-1/2}\mathcal A,\qquad
 S=\bar{\mathcal A}\bar{\mathcal A}^*=D_p^{-1/2}LD_p^{-1/2},\qquad
 \Omega=\bar{\mathcal A}^*\bar{\mathcal A} .
\]
Since $L\succeq0$ and $S=I-D_p^{-1/2}(P\circ P)D_p^{-1/2}$ with $P\circ P\succeq0$
(Schur product theorem),
\begin{equation}\label{eq:S}
 0\preceq S\preceq I,\qquad 0\preceq \Omega\preceq I,\qquad
 \tr(I-S)=\sum_i\frac{P_{ii}^2}{p_i}=d .
\end{equation}
Fix an orthonormal eigenbasis of $S$, and for each of its members $y$ with
eigenvalue $\mu>0$ put $f_y=D_p^{-1/2}y$ and
$Z_y=\mu^{-1/2}\bar{\mathcal A}^*y=\mu^{-1/2}N_{f_y}$. Since $\Omega Z_y=\mu Z_y$ and
$\ip{Z_y}{Z_{y'}}=\ip y{Sy'}/\sqrt{\mu\mu'}$, the $Z_y$ form an orthonormal
eigenbasis of $\Omega$ on $(\ker \Omega)^\perp$. Sums $\sum_{\mu>0}$ below run over these
basis members (with multiplicity). Since $p_i\ge a/2$,
\begin{equation}\label{eq:finfty}
 \norm{f_y}_\infty\le\sqrt{2/a}.
\end{equation}

\begin{definition}[Filtered covariance]
For $\rho>0$ let $C_\rho=\rho(I-\Omega)(\rho I+\Omega)^{-1}$ on $\mathcal Z$, let
$Z\sim N(0,C_\rho/n)$, and $Y=Y_Z$.
\end{definition}

\begin{lemma}\label{lem:filter}
\begin{enumerate}[label=(\alph*)]
\item $0\preceq C_\rho\preceq I$, and $C_\rho=I$ on $\ker \Omega$.
\item $\Cov(\mathcal AZ)=n^{-1}D_p^{1/2}\,\rho S(I-S)(\rho I+S)^{-1}D_p^{1/2}\preceq
(\rho/n)D_p$.
\item $I-C_\rho=\Omega+R_\rho$ with $R_\rho=\sum_{\mu>0}r_\mu\,Z_y\otimes Z_y$,
$r_\mu=\mu(1-\mu)/(\rho+\mu)\ge0$, and
\[
 \sum_{\mu>0}r_\mu\le d,\qquad \sum_{\mu>0}\frac{r_\mu}{\sqrt\mu}\le\frac d{2\sqrt\rho},
 \qquad\sum_{\mu>0}\frac{r_\mu}\mu\le\frac d\rho .
\]
\end{enumerate}
\end{lemma}

\begin{proof}
(a) The eigenvalues of $C_\rho$ are $\rho(1-\mu)/(\rho+\mu)\in[0,1]$. (b) By
singular value calculus $\bar{\mathcal A}C_\rho\bar{\mathcal A}^*=\rho S(I-S)(\rho
I+S)^{-1}$, whose eigenvalues $\rho\mu(1-\mu)/(\rho+\mu)$ are at most $\rho$. (c)
On the $\mu$-eigenspace, $1-\rho(1-\mu)/(\rho+\mu)=\mu+\mu(1-\mu)/(\rho+\mu)$. The
sums follow from $\mu/(\rho+\mu)\le1$, $\sqrt\mu/(\rho+\mu)\le1/(2\sqrt\rho)$,
$1/(\rho+\mu)\le1/\rho$, and $\sum_{\mu>0}(1-\mu)\le\tr(I-S)=d$.
\end{proof}

\begin{remark}
Only these properties of the filter $\phi(\mu)=\rho(1-\mu)/(\rho+\mu)$ are used:
$\phi(0)=1$ (so $C_\rho=I$ on $\ker \Omega$), $0\le\phi(\mu)\le1-\mu$,
$\mu\phi(\mu)\le\rho$, and the sums in (c). The hard cutoff $(1-\mu)\one[\mu<\rho]$
satisfies them with $\sum r_\mu/\sqrt\mu\le d/\sqrt\rho$, which doubles some
constants below; the rational filter keeps every object a rational function of
$S$, which is convenient for formalisation.
\end{remark}

\subsection{Commutator identities}

For $Z\in\mathcal Z$ define the quadratic diagonal and its polarisation
\[
 q_i(Z)=\norm{Z^Tv_i}^2-\norm{Zu_i}^2,\qquad
 q_i(Z,M)=\ip{Z^Tv_i}{M^Tv_i}-\ip{Zu_i}{Mu_i},
\]
so that $q(Z+M)=q(Z)+2q(Z,M)+q(M)$ and, by Cauchy--Schwarz,
\begin{equation}\label{eq:q-cs}
 |q_i(Z,M)|\le\norm{(Y_Z)_i}\,\norm{(Y_M)_i},\qquad |q_i(M)|\le\norm{(Y_M)_i}^2 .
\end{equation}

\begin{lemma}\label{lem:commutator}
Let $f\in\R^n$, $F=\Diag(f)$, and $R=I-2P$ (an orthogonal matrix). Then:
\begin{enumerate}[label=(\alph*)]
\item $q(N_f)=\mathcal A\Gamma_f$ with $\Gamma_f=V^TFRFU$, and
$\fro{\Gamma_f}\le2\norm f_\infty\fro{N_f}$;
\item $T_f:=Y_{N_f}=FP+PF-2PFP$ satisfies
$[X,T_f]=RF[X,P]+[X,P]FR$ for diagonal $X$; hence
$\Lap(T_f)\preceq4\norm f_\infty^2L$, and
$x^T\Lap(T_{e_k})x\le2\sum_jP_{kj}^2(x_k-x_j)^2$.
\end{enumerate}
\end{lemma}

\begin{proof}
(a) For $x\in\R^n$ put $X=\Diag(x)$, $B_x=V^TXV$, $A_x=U^TXU$. Then
$\ip x{q(Z)}=\ip Z{B_xZ-ZA_x}$ for every $Z$, and, since $X$ and $F$ commute,
\[
 B_xN_f-N_fA_x=V^T[X(I-P)F-FPX]U=V^T[F(I-P)X-XPF]U=B_fN_x-N_xA_f .
\]
Hence $\ip x{q(N_f)}=\ip{N_f}{B_fN_x-N_xA_f}=\ip{B_fN_f-N_fA_f}{N_x}$. As
$\ip M{N_x}=\sum_ix_iv_i^TMu_i=\ip{\mathcal AM}x$ for every $M$, this says
$q(N_f)=\mathcal A(B_fN_f-N_fA_f)$, and
$B_fN_f-N_fA_f=V^TF(I-P)FU-V^TFPFU=\Gamma_f$. The norm bound follows from
$\opn{B_f},\opn{A_f}\le\norm f_\infty$.

(b) $Y_{N_f}=VV^TFUU^T+UU^TFVV^T=(I-P)FP+PF(I-P)$. Since $X$ commutes with $F$,
$[X,FP]=F[X,P]$, $[X,PF]=[X,P]F$ and $[X,PFP]=[X,P]FP+PF[X,P]$, which gives the
commutator formula. Then $\fro{[X,T_f]}\le2\norm f_\infty\fro{[X,P]}$, i.e.\ the
first bound (\cref{lem:energy}). For $f=e_k$ and $E_k=\Diag(e_k)$, $E_k[X,P]$ and $[X,P]E_k$ are
the $k$th row and column of $[X,P]$, each of squared norm
$\sum_jP_{kj}^2(x_k-x_j)^2$.
\end{proof}

\subsection{The second-order mean and the drift}

\begin{lemma}[Mean]\label{lem:mean}
$\E q(Z)=a\one-p-b_0-b_*$ with $b_*:=\mathcal Am_*$, where
\[
 b_0=\frac1nL(p^{-1}),\quad \norm{b_0}_\infty\le\frac{12\eps}n,\qquad
 m_*=\frac1n\sum_{\mu>0}\frac{r_\mu}\mu\,\Gamma_{f_y},\quad
 \fro{m_*}^2\le\frac{2a}\rho .
\]
\end{lemma}

\begin{proof}
For $G\sim N(0,I)$ on $\mathcal Z$, $\E\norm{G^Tv_i}^2=d(1-p_i)$ and
$\E\norm{Gu_i}^2=(n-d)p_i$, so the unfiltered covariance $I/n$ contributes
$a\one-p$. The covariance $\Omega=\sum_j\zeta_j\otimes\zeta_j$, $\zeta_j=\bar{\mathcal
A}^*e_j=v_ju_j^T/\sqrt{p_j}$, contributes
\[
 \frac1n\sum_jq_i(\zeta_j)=\frac1n\sum_j\frac{(I-P)_{ij}^2p_j-(1-p_j)P_{ij}^2}{p_j}
 =\frac1n\Bigl(1-\sum_j\frac{P_{ij}^2}{p_j}\Bigr)=\frac1n(Lp^{-1})_i ,
\]
using $\sum_j(I-P)_{ij}^2=1-p_i$ and $\sum_jP_{ij}^2=p_i$. Here
$(Lp^{-1})_i=\sum_jP_{ij}^2(p_i^{-1}-p_j^{-1})$ and
$|p_i^{-1}-p_j^{-1}|\le2\eps a/((1-\eps)a)^2\le8\eps/a$, so
$|(Lp^{-1})_i|\le12\eps$. Finally $R_\rho/n$ contributes
$\frac1n\sum_\mu r_\mu q(Z_y)=\frac1n\sum_\mu\frac{r_\mu}\mu q(N_{f_y})=\mathcal
Am_*$ by \cref{lem:commutator}(a), and with $\fro{N_{f_y}}=\sqrt\mu$,
\eqref{eq:finfty} and \cref{lem:filter}(c),
\[
 \fro{m_*}\le\frac2n\sqrt{\frac2a}\sum_{\mu>0}\frac{r_\mu}{\sqrt\mu}
 \le\frac2n\sqrt{\frac2a}\cdot\frac d{2\sqrt\rho}=\sqrt{\frac{2a}\rho}.\qedhere
\]
\end{proof}

\begin{remark}[Why a drift]\label{rem:drift}
The filter removes covariance from the directions $Z_y$, and with it their share
$t^2b_*=t^2\mathcal Am_*$ of the second-order drift towards $a\one$. This bias
need not be below the tolerance $\sigma/(12\Theta)$ used later: for instance, for
many small blocks joined through a common row it is of order $at^2$ at that row.
(Since $\Gamma_f=V^T[(ff^T)\circ(I-2P)]U$ is quadratic in $f$, and kernel modes
$Sy=0$ have $N_{f_y}=0$ and hence $\Gamma_{f_y}=0$, the spectral identity
$\sum_y\frac{1-\mu}{\rho+\mu}yy^T=(I-S)(\rho I+S)^{-1}$ shows that $m_*$ does not
depend on the choice of eigenbasis:
$m_*=\frac1nV^T\bigl[(D_p^{-1/2}(I-S)(\rho I+S)^{-1}D_p^{-1/2})\circ(I-2P)\bigr]U$.)
However, $t^2b_*$ is exactly the \emph{first-order} diagonal change produced
by the deterministic horizontal direction $\frac t2m_*$, whose norm is
$\frac t2\fro{m_*}\le\sqrt{a/2}/D$. Adding this direction to the noise cancels
the bias exactly, at the price of cross terms of size $O(\sigma/D)$.
\end{remark}

\subsection{The expected graph}

Since $\E Y_{ij}^2$ is linear in the covariance, $\E_\Sigma\Lap(Y)=\sum_k\lambda_k
\Lap(Y_{\psi_k})$ for every representation $\Sigma=\sum_k\lambda_k\psi_k\otimes\psi_k$
with real $\lambda_k$; we use it for $C_\rho/n=(I-\Omega-R_\rho)/n$ with
$\Omega=\sum_j\zeta_j\otimes\zeta_j$.

\begin{lemma}\label{lem:expected-graph}
$\displaystyle\E\Lap(Y)\succeq\frac a{4n}\mathsf K-\Bigl(\frac2n+\frac8d+\frac8\rho\Bigr)L$.
\end{lemma}

\begin{proof}
\emph{Unfiltered noise.} For $Z_0\sim N(0,I/n)$ and $i\ne j$,
$(Y_{Z_0})_{ij}=v_i^TZ_0u_j+v_j^TZ_0u_i$ has variance
$n^{-1}\fro{v_iu_j^T+v_ju_i^T}^2=n^{-1}(p_i+p_j-2p_ip_j-2P_{ij}^2)$. As
$p\le\frac34$, $p_i+p_j-2p_ip_j\ge\frac14(p_i+p_j)\ge\frac a4$, so
$\E\Lap(Y_{Z_0})\succeq\frac a{4n}\mathsf K-\frac2nL$.

\emph{Base loss.} $Y_{\zeta_k}=T_{e_k}/\sqrt{p_k}$, so by
\cref{lem:commutator}(b) and $p_k\ge a/2$,
$\sum_kx^T\Lap(Y_{\zeta_k})x\le\frac2a\cdot2\sum_{k,j}P_{kj}^2(x_k-x_j)^2
=\frac8a\,x^TLx$; the base part $\Omega/n$ therefore removes at most $\frac8dL$.

\emph{Residual loss.} $\Lap(Y_{Z_y})=\mu^{-1}\Lap(T_{f_y})\preceq\frac8{a\mu}L$ by
\cref{lem:commutator}(b) and \eqref{eq:finfty}, so $R_\rho/n$ removes at most
$\frac8{an}\sum_\mu\frac{r_\mu}\mu L\le\frac8\rho L$.
\end{proof}

\subsection{A good sample}

Let $\ell_i=\frac1n\sum_\mu r_\mu\norm{(Y_{Z_y})_i}^2$ be the variance that $R_\rho/n$
removes from row $i$. Since $\fro{Y_{Z_y}}^2=2$, $\sum_i\ell_i\le\frac2n\sum_\mu
r_\mu\le2a$. Fix $\delta_0=1/6400$ and let
\[
 \mathfrak B=\{i:\ell_i>\delta_0a\},\qquad b=|\mathfrak B|\le2/\delta_0=:b_{\max},
\]
a deterministic set depending only on $U$ and $\rho$. Write
$\mathcal C(Y)=\sum_{i<j}2P_{ij}Y_{ij}(e_i-e_j)(e_i-e_j)^T$, so that
$\Lap(P+tY)=L+t\mathcal C(Y)+t^2\Lap(Y)$.

\begin{lemma}\label{lem:sample}
There is an absolute $h\in(0,1)$ and, for each $\eta'\in(0,\frac1{32}]$, a
constant $A_{\eta'}$ such that the following holds for all $\rho>0$ and
$0<t\le1$ whenever $d\ge A_{\eta'}\log(2n)$. With positive probability:
\begin{enumerate}[label=\textup{(S\arabic*)}]
\item\label{S1} $\opn Z\le16$, $\max_i\norm{Z^Tv_i}^2\le3a$ and
$\max_i\norm{Y_i}^2\le800a$;
\item\label{S2} $\max_i|(\mathcal AZ)_i|\le3\sqrt{\rho a\log(2n)/n}$ and
$\max_i|q_i(Z)-\E q_i(Z)|\le\eta'a$;
\item\label{S3} every $i\notin\mathfrak B$ has at least $0.95n$ indices
$j\ne i$ with $|P_{ij}+tY_{ij}|\ge h\,t\sqrt{a/n}$;
\item\label{S4} for every $x$ supported on $\mathfrak B$,
\[
 |t\,x^T\mathcal C(Y)x|\le\eta'\bigl(x^TLx+at^2\norm x^2\bigr),\qquad
 |x^T(\Lap(Y)-\E\Lap(Y))x|\le\eta'a\norm x^2 .
\]
\end{enumerate}
\end{lemma}

\begin{proof}
Since $C_\rho\preceq I$, every linear functional $\ip Z\Psi$ has variance at
most $\fro\Psi^2/n$; we use the tools of \cref{app:gauss}.

\ref{S1}: The net bound (\cref{lem:gauss}(c)) with $\sigma^2=1/n$ gives
$\Prob(\opn Z>16)\le2\cdot5^ne^{-8n}$. The vector $Z^Tv_i$ has covariance
$\preceq n^{-1}I_d$ and trace $\le a$; \cref{lem:gauss}(a) with $u=2a$ (so
$\min\{u^2/\fro\cdot^2,u/\opn\cdot\}\ge2d$) and a union bound give
$\max\norm{Z^Tv_i}^2\le3a$ off probability $2ne^{-cd}$. Then
$\norm{Y_i}\le\norm{Z^Tv_i}+\norm{Zu_i}\le\sqrt{3a}+16\sqrt{3a/2}$.

\ref{S2}: $(\mathcal AZ)_i$ is Gaussian with variance $\le\rho p_i/n\le3\rho
a/(2n)$ (\cref{lem:filter}(b)); a union bound gives the first claim off
probability $2n(2n)^{-3}$. The form $q_i(Z)=\ip Z{\mathsf M_iZ}$,
$\mathsf M_iZ=v_iv_i^TZ-Zu_iu_i^T$, has $\opn{\mathsf M_i}\le1$ and
$\fro{\mathsf M_i}^2\le2(d+np_i^2)\le7d$; after the covariance factor
$(C_\rho/n)^{1/2}$ these become $\le1/n$ and $\le\sqrt{7d}/n$, and
\cref{lem:gauss}(a) with $u=\eta'a$ gives failure probability $2n\exp(-c\eta'^2d)$.

\ref{S3}: \emph{Variance bookkeeping.} By the proof of
\cref{lem:expected-graph}, for $i\ne k$ with $P_{ik}^2\le a/16$ the unfiltered
variance of $Y_{ik}$ is at least $a/(8n)$; at most $24$ indices $k$ have
$P_{ik}^2>a/16$. Row $i$ of $T_{e_j}$ is
$\delta_{ij}e_j^TP+P_{ij}e_j^T-2P_{ij}e_j^TP$, of squared norm at most
$3(\delta_{ij}p_j+P_{ij}^2+4P_{ij}^2p_j)$, so the base part removes from row $i$
total variance $\frac1n\sum_j\norm{e_i^TT_{e_j}}^2/p_j\le\frac6{an}\cdot5p_i\le
\frac{45}n$, which exceeds $\frac a{32n}$ at no more than $1440n/d$ positions.
For $i\notin\mathfrak B$ the residual part removes at most $\delta_0a$ from row
$i$, hence more than $\frac a{32n}$ at no more than $32\delta_0n$ positions. So for
$d\ge10^6$ (hence $n\ge2\cdot10^6$, and the excluded indices number at most
$25+0.0065n\le0.01n$) every $i\notin\mathfrak B$ has at least $0.99n$ indices $k\ne i$
with $\Var(Y_{ik})\ge\frac a{16n}$.

\emph{Small ball.} Let $\phi:\R\to[0,1]$ be $C^1$, equal to $1$ on $[-h,h]$,
vanishing outside $[-2h,2h]$, with $|\phi'|\le2/h$, and for $i\notin\mathfrak B$
let $F_i=\frac1n\sum_{k\ne i}\phi\bigl((P_{ik}+tY_{ik})/(t\sqrt{a/n})\bigr)$. By
the density bound for a Gaussian of variance $\ge\frac{at^2}{16n}$,
$\E F_i\le0.01+\frac{4h}{\sqrt{2\pi/16}}\le0.01+7h$. As a function of the row
$(Y_{ik})_k$, $F_i$ has gradient norm at most
$\frac1n\cdot\frac2h\cdot\sqrt{n/a}\cdot\sqrt n=2/(h\sqrt a)$; the row depends on
$Z$ through a linear map of norm $\le2$, and $Z=(C_\rho/n)^{1/2}g$ with $g$
standard. So $F_i$ is $4/(h\sqrt d)$-Lipschitz in $g$, and \cref{lem:gauss}(b)
gives $\Prob(F_i>\E F_i+0.01)\le2\exp(-c\,h^2d)$. Take $h=1/500$, so that
$\E F_i+0.01\le0.04$, and a union bound over $i$. On the complement, at most
$0.04n$ indices $k\ne i$ have $|P_{ik}+tY_{ik}|<ht\sqrt{a/n}$, so at least
$n-1-0.04n\ge0.95n$ have the reverse inequality.

\ref{S4}: Only second moments are needed, because $\mathfrak B$ is deterministic
and bounded. Let $\Pi_{\mathfrak B}$ be the coordinate projection onto
$\R^{\mathfrak B}$, $J_{\mathfrak B}=L_{\mathfrak B\mathfrak B}+at^2I_{\mathfrak B}$
and $\omega_{ij}=J_{\mathfrak B}^{-1/2}\Pi_{\mathfrak B}(e_i-e_j)$, so
$\norm{\omega_{ij}}^2\le2/(at^2)$ and
$\sum_{i<j}P_{ij}^2\omega_{ij}\omega_{ij}^T=J_{\mathfrak B}^{-1/2}L_{\mathfrak
B\mathfrak B}J_{\mathfrak B}^{-1/2}\preceq I$. For $x$ supported on $\mathfrak
B$, $tx^T\mathcal C(Y)x=\ip{J_{\mathfrak B}^{1/2}x}{\Xi J_{\mathfrak B}^{1/2}x}$
with $\Xi=\sum_{i<j}2tP_{ij}Y_{ij}\,\omega_{ij}\omega_{ij}^T$, and
$x^TJ_{\mathfrak B}x=x^TLx+at^2\norm x^2$. The coefficients $(Y_{ij})_{i<j}$
have covariance $\preceq\frac2nI$ (the map $Z\mapsto Y_Z$ is $\sqrt2$ times an
isometry), so
\[
 \E\fro\Xi^2\le\frac{2t^2}n\sum_{i<j}4P_{ij}^2\norm{\omega_{ij}}^4
 \le\frac{16}{an}\tr\bigl(J_{\mathfrak B}^{-1/2}L_{\mathfrak B\mathfrak B}
 J_{\mathfrak B}^{-1/2}\bigr)\le\frac{16b}d ,
\]
and Markov's inequality gives the first half off probability $16b/(d\eta'^2)$.
The principal $\mathfrak B$-block of $\Lap(Y)-\E\Lap(Y)$ has diagonal entries
$\sum_{j\ne i}(Y_{ij}^2-\E Y_{ij}^2)$, a centred Gaussian quadratic form of
variance at most $2\cdot\frac2n\cdot\E\norm{Y_i}^2\le\frac{12a}n$
(\cref{lem:gauss}(a), with $\E\norm{Y_i}^2\le a(1-p_i)+(1-a)p_i\le3a$), and
off-diagonal entries $-(Y_{ik}^2-\E Y_{ik}^2)$ of variance at most $8/n^2$. Its
expected squared Frobenius norm is at most $20b^2a/n=20b^2a^2/d$, and Markov
gives the second half off probability $20b^2/(d\eta'^2)$.

Finally choose $A_{\eta'}$ so large that $d\ge A_{\eta'}\log(2n)$ implies
$d\ge10^6$ and all failure probabilities above are at most $\frac1{10}$; their
sum is less than $1$.
\end{proof}

\subsection{The drifted retraction}

Fix $D\ge12$ and $t\in(0,1]$, set $\rho=D^2t^2$, take a sample $Z$ as in
\cref{lem:sample}, and let $m_*$ be as in \cref{lem:mean}. Define
\[
 W=Z+\tfrac t2m_*,\qquad G=I+t^2W^TW,\qquad U_t=(U+tVW)G^{-1/2},\qquad
 P_t=U_tU_t^T .
\]
Since $U^TV=0$, $(U+tVW)^T(U+tVW)=G$ and $U_t$ is Parseval.

\begin{lemma}\label{lem:retraction}
There is an absolute $C_E$ (depending only on the bounds in \ref{S1} and on
$D\ge12$) such that:
\begin{enumerate}[label=(\alph*)]
\item $\fro{U_t-U}^2\le C_Et^2d$;
\item every row of $\mathsf E=P_t-P-tY$ satisfies
$\sum_j\mathsf E_{ij}^2\le at^2(D^{-2}+C_Et^2)$;
\item with $\norm{\mathsf r}_\infty\le C_Eat^3$,
\[
 \diag P_t-a\one=(1-t^2)(p-a\one)+2t\mathcal AZ-t^2b_0+t^2\bigl(q(Z)-\E q(Z)\bigr)
 +t^3q(Z,m_*)+\tfrac{t^4}4q(m_*)+\mathsf r ;
\]
\item $\norm{\diag P_t-a\one}_\infty\le at^2\bigl(\frac{\eps}{t^2}
+6D\sqrt{\log(2n)/d}+\frac{12\eps}d+\eta'+\frac{40}D+\frac1{2D^2}+C_Et\bigr)$.
\end{enumerate}
\end{lemma}

\begin{proof}
By \cref{lem:mean}, $\frac t2\opn{m_*}\le\frac t2\fro{m_*}\le\sqrt{a/2}/D\le\frac1{24}$,
so with \ref{S1}: $\opn W\le17$, $\fro W\le17\sqrt d$,
$\norm{W^Tv_i}\le2\sqrt a$, $\norm{Wu_i}\le17\sqrt{3a/2}$, and
$\norm{(Y_{m_*})_i}^2\le\opn{m_*}^2\le2a/\rho$. Hence $\opn{G^{-1}-I}$ and
$\opn{G^{-1/2}-I}$ are at most $289t^2$, $\opn{U+tVW}\le18$ and
$\norm{u_i+tW^Tv_i}\le4\sqrt a$.

(a) $\fro{U_t-U}\le\fro{U+tVW}\opn{G^{-1/2}-I}+t\fro W\le Ct\sqrt d$.

(b) $\mathsf E=\frac{t^2}2Y_{m_*}+\mathsf E'$ with
$\mathsf E'=P_t-P-tY_W=(U+tVW)(G^{-1}-I)(U+tVW)^T+t^2VWW^TV^T$, whose rows have
norm at most $Ct^2\sqrt a$. The rows of $\frac{t^2}2Y_{m_*}$ have squared norm at
most $\frac{t^4}4\cdot\frac{2a}\rho=\frac{at^2}{2D^2}$. Use $(\alpha+\beta)^2\le
2\alpha^2+2\beta^2$.

(c) Let $\xi_i=u_i+tW^Tv_i$, the $i$th row of $U+tVW$. Inserting
$G^{-1}=I-t^2W^TW+t^4(W^TW)^2G^{-1}$ into $(P_t)_{ii}=\xi_i^TG^{-1}\xi_i$ and using
$u_i^TW^Tv_i=(\mathcal AW)_i$,
\[
 (P_t)_{ii}=p_i+2t(\mathcal AW)_i+t^2q_i(W)-2t^3\ip{Wu_i}{WW^Tv_i}
 -t^4\norm{WW^Tv_i}^2+t^4\xi_i^T(W^TW)^2G^{-1}\xi_i ,
\]
and the last three terms are $O(at^3)$. Now $2t\mathcal AW=2t\mathcal AZ+t^2\mathcal
Am_*$, $q(W)=q(Z)+tq(Z,m_*)+\frac{t^2}4q(m_*)$, and
$\E q(Z)=a\one-p-b_0-\mathcal Am_*$ by \cref{lem:mean}; the terms $\pm t^2\mathcal
Am_*$ cancel.

(d) In (c): $(1-t^2)\norm{p-a\one}_\infty\le\eps a$;
$2t\norm{\mathcal AZ}_\infty\le6tD t\sqrt{a\log(2n)/n}$ by \ref{S2}, which is
$6Dat^2\sqrt{\log(2n)/d}$; $t^2\norm{b_0}_\infty\le12\eps t^2/n$; the fluctuation is
at most $\eta'at^2$; by~\eqref{eq:q-cs} and \ref{S1},
$t^3|q_i(Z,m_*)|\le t^3\sqrt{800a}\sqrt{2a/\rho}=40at^2/D$ and
$\frac{t^4}4|q_i(m_*)|\le\frac{t^4}4\cdot\frac{2a}\rho=\frac{at^2}{2D^2}$.
\end{proof}

\begin{lemma}[Graph of the seed]\label{lem:seed-graph}
There are absolute $t_*\in(0,1]$ and $C_{\rm core}$ such that the following
holds. Let $\eta'\le\frac1{32}$, $d\ge A_{\eta'}\log(2n)$, $D\ge13/h$,
$b\le n/10$ and $t\le t_*$, and let $Z$ satisfy \ref{S1}--\ref{S4}. Then
$H(L_{P_t})\le C_{\rm core}/(at^2)$.
\end{lemma}

\begin{proof}
\emph{Core.} In a row $i\notin\mathfrak B$, by \cref{lem:retraction}(b), at most
$at^2(D^{-2}+C_Et^2)\big/\bigl(\frac{h^2}4\frac{at^2}n\bigr)\le0.05n$ indices have
$|\mathsf E_{ij}|\ge\frac h2t\sqrt{a/n}$, once $D\ge13/h$ and
$t^2\le h^2/(160C_E)$. With \ref{S3}, every $i\notin\mathfrak B$ has at least
$0.9n$ indices $j\ne i$ with $(P_t)_{ij}^2\ge\frac{h^2}4\frac{at^2}n$.

\emph{Block.} For $x$ supported on $\mathfrak B$, \cref{lem:expected-graph},
$x^T\mathsf Kx\ge(n-b)\norm x^2$ and \ref{S4} give
\[
 x^T\Lap(P+tY)x\ge\Bigl(1-\eta'-t^2\bigl(\tfrac2n+\tfrac8d\bigr)-\tfrac8{D^2}\Bigr)x^TLx
 +at^2\Bigl(\tfrac14\bigl(1-\tfrac bn\bigr)-2\eta'\Bigr)\norm x^2
 \ge\frac{at^2}8\norm x^2 ,
\]
since the first coefficient is positive and $\frac14\cdot\frac9{10}-\frac1{16}\ge\frac18$.
By \eqref{eq:lap-basic} and \cref{lem:retraction}(b),
$L_{P_t}=\Lap(P+tY+\mathsf E)\succeq\frac12\Lap(P+tY)-2at^2(D^{-2}+C_Et^2)I$,
which is at least $\frac{at^2}{32}$ on vectors supported on $\mathfrak B$ once
$t^2\le1/(128C_E)$.

\emph{Barrier.} \Cref{lem:core} with $\mathcal B=\mathfrak B$, part (i) with
$\vartheta=0.4$, $\gamma=h^2at^2/4$, and part (ii) with $\lambda=at^2/32$, gives
$H(L_{P_t})\le\frac{1+b}{0.1h^2at^2}+\frac{32\sqrt b}{at^2}$. Put
$C_{\rm core}=10(1+b_{\max})/h^2+32\sqrt{b_{\max}}$.
\end{proof}

\subsection{Proof of \texorpdfstring{\cref{thm:moderate}}{Theorem}}

\emph{Constants.} The numbers $\delta_0,b_{\max},h,C_E,t_*,C_{\rm core}$ are
absolute and do not depend on $D$ (for $D\ge12$). Put
$\Theta=\max\{2,C_{\rm core},C_E\}$ and choose, in this order: $D\ge\max\{13/h,
500\Theta\}$, so that $\frac{40}D+\frac1{2D^2}\le\frac1{12\Theta}$; $K_0=12\Theta$
and $\eta'=\min\{\frac1{32},\frac1{12\Theta}\}$; $A\ge A_{\eta'}$ so large that
$d\ge A\log(2n)$ implies $6D\sqrt{\log(2n)/d}\le\frac1{12\Theta}$ and
$b_{\max}\le n/10$; and $\eps_0\le\frac12$ so that $t=\sqrt{K_0\eps}$ satisfies
$t\le t_*$, $C_Et\le\frac1{12\Theta}$ and $12\eps\le\frac1{12\Theta}$ for
$\eps\le\eps_0$. None of these choices depends on $n,d,\eps$ or $U$.

\emph{Seed.} With $t^2=K_0\eps$, \cref{lem:retraction}(d) bounds the diagonal
error of $P_t$ by $at^2\cdot\frac6{12\Theta}=\frac{at^2}{2\Theta}$;
\cref{lem:seed-graph} bounds the barrier constant by $\Theta/(at^2)$; and
\cref{lem:retraction}(a) bounds $\fro{U_t-U}^2$ by $\Theta t^2d$. So $U_t$ is a
$(\Theta,t)$-seed for $U$, and \cref{cor:seed} gives an ENP frame $\widehat U$
with $\fro{\widehat U-U}^2\le3\Theta t^2d=36\Theta^2\eps d$.
\qed

\begin{remark}
The drift keeps the spectral input local. The bias $t^2b_*$ could also be
removed after the perturbation, for instance by a second scaling; since
$\ip x{b_*}^2\le\frac{2a}\rho x^TLx$ this is possible, but it requires two-sided
control of $L_{P_t}$ on all of $\one^\perp$. With the drift, the only spectral
input is the gap on the bounded deterministic block $\mathfrak B$, which needs
second moments only.
\end{remark}
